\section{Discussion}
\label{sec:discussion}

\subsection{What the transcriptome does}

The induced set reads as a single coherent strategy rather than a collection of
stress responses, and it agrees in direction with earlier work on shifting
nitrogen supply in a related diatom \citep{alexander2015metabolic}. NRT2.1 and AMT1.3 raise the affinity of the cell for the two
inorganic nitrogen species still present at trace concentration, glutamine
synthetase III raises the rate at which whatever is taken up is fixed into
glutamine, and the urea and purine catabolism group reclaims nitrogen already
committed to nucleotides. The order of effect sizes matches that logic: uptake
first, assimilation second, internal salvage third.

The repressed set is the mirror image and is dominated by two groups that both
represent large nitrogen investments. Fucoxanthin chlorophyll proteins are
among the most abundant proteins in a diatom cell, and cytosolic ribosomes are
the largest single sink for cellular nitrogen. Repressing both frees nitrogen
without requiring any new machinery. The drop in
$F_v/F_m$ from 0.66 to 0.38 in Table~\ref{tab:physiology} is consistent with
the antenna repression being realised at the protein level and not only in
transcript abundance.

The urea cycle result deserves a note of caution. Carbamoyl phosphate synthase
and the urea transporter DUR3 both clear our threshold, but their adjusted
$p$ values are three to four orders of magnitude weaker than those of the
transporters, and in the leave-one-out reanalysis DUR3 fell below significance
when NL-3 was dropped. A study designed to resolve the urea cycle specifically
would need more replicates rather than more depth, because the limiting
quantity is between-flask variance and not counting noise.

\subsection{Limitations}

Three limitations bound what this experiment can claim. The design contrasts a
single time point, so it cannot separate the immediate response to nitrogen
depletion from the acclimated state reached after several generations. The
comparison is between regimes that also differ in growth rate, and a growth
rate difference alone produces ribosomal protein repression, so part of the
translation signal is confounded. Finally, transcript abundance is a poor proxy
for flux, and the salvage pathways in particular are known to be regulated
post-translationally in this species.

\subsection{Recommendations for the next cohort}

Retain the duplication threshold at 15 percent as a flag rather than a
rejection criterion, and require the leave-one-out check that we ran here
before any flagged library is kept. Add a mid-limitation time point at 48 hours
so that the response can be ordered in time. Increase biological replication
from three to five flasks per regime at the same total sequencing budget: the
analysis in Section~\ref{sec:results} shows that the significant set is
variance-limited rather than depth-limited beyond about 25 million read pairs.

\section*{Data availability}

Raw reads, the container manifest, the design matrix, and the analysis scripts
are archived in the unit repository under experiment series NL-08. Count
matrices and the full differential expression table are included with this
report.

\section*{Contributions}

I. Karvonen maintained the cultures and performed the extractions.
A. Diallo prepared the libraries and ran the sequencing.
R. Ferreira-Bax wrote the analysis pipeline and the statistical model. All
three authors reviewed the quality control decisions described in
Section~\ref{sec:results}.
